\section{课程导论与2025年软件开发现状}

\subsection{为什么需要这门课}

2025年，软件开发正经历一场深刻的范式转变。以 Windsurf、Cursor、Claude Code 等为代表的 AI 编程工具正在以前所未有的速度改变开发者的工作方式。Mihail Eric 在课程开篇指出了两个关键信号：

\begin{importantbox}{核心观点：AI 不会取代开发者}
``You won't be replaced by AI. You'll be replaced by a competent engineer who knows how to use AI.''

AI 本身不会取代软件工程师，但善于使用 AI 的工程师将取代不善于使用 AI 的工程师。这门课的目标是让学生成为前者。
\end{importantbox}

好消息是：软件开发者拥有历史上最高的生产力潜力。借助 AI 编程工具，工程师可以以前所未有的速度学习新技术栈和工具。

\begin{warningbox}{这不是 ``Vibe Coding'' 课}
课程明确强调：这不是一门教你用 AI 随意生成代码的课程。重点在于理解 AI 编程工具的原理、掌握高质量的 human-agent 协作模式，而非盲目依赖 AI。
\end{warningbox}

\subsection{课程核心理念}

课程围绕以下关键理念展开：

\begin{enumerate}
    \item \textbf{Human-Agent Engineering}：专注于 AI 尚未替代的技能——业务理解、技术架构决策、成为 Tech Lead
    \item \textbf{LLM 的能力取决于你}：好的 context 产生好的代码；如果你自己都无法理解你的代码库，LLM 也不能
    \item \textbf{大量阅读和审查代码}：学会辨别好代码与坏代码，培养``品味''（taste）
    \item \textbf{激进地实验}：目前没有成熟的软件开发范式，每个人都在摸索中
\end{enumerate}

\subsection{课程安排概览}

10周课程涵盖 AI 编程工具链的各个层面：从 LLM 基础、编程 Agent 构建、AI IDE、终端工具、测试与安全、代码审查、UI 自动化构建、DevOps，到 AI 软件工程的未来展望。每周邀请行业顶尖创始人做客座讲座。

\subsection{本章小结}

本节确立了课程的基调：AI 正在改变软件开发，但核心竞争力仍然在于工程师自身的理解力、判断力和对工具的熟练运用。关键是学会如何与 AI Agent 高效协作，而非被 AI 所取代。

